\section{Discussion}
\label{sec:discussion}

\subsection{Privacy Regulation as a Valuation Shock}

The results show that the 2021 privacy regime changed the valuation of Chinese digital firms. The average effect is negative for high-TDDI firms, but the distribution of the effect is uneven. B2B data-service firms lose value, while consumer internet firms with first-party data do not. This pattern is consistent with the view that privacy regulation redistributes rents across business models rather than simply reducing the total amount of data in the economy.

The comparison with prior work is instructive. Studies of the GDPR document average effects on web traffic, advertising, and firm performance \cite{goldberg2024,aridor2020}. Studies of ATT document effects on e-commerce, app monetization, and the app ecosystem \cite{aridor2024,kesler2023,cheyre2024}. These designs compare regulated and unregulated markets, or exposed and unexposed users. Our design compares firms within the same regulated market and uses predetermined text-based exposure. The results show that the cross-sectional distribution of exposure is economically important: average effects can mask large losses in data-service segments and negligible effects in first-party-data segments.

The B2B--B2C contrast also speaks to the JTAER literature on privacy regulation and platform competition. Recent work shows that privacy regulation affects retailer resilience through customer-acquisition dependence and traffic architecture \cite{chen2026}. Our results complement this finding by showing that the valuation effect is concentrated in firms that sell data-intensive services to other firms. Privacy regulation operates not only through consumer consent but also through the business-to-business data supply chain.

\subsection{Implications for Data-Asset Accounting}

Data-asset accounting currently focuses on the recognition and measurement of internally generated data resources. The results suggest that third-party data dependence is a distinct risk factor. Two firms with the same book assets and the same reported revenue can have different exposure to privacy regulation because their data inputs differ. A TDDI-style measure can be constructed from public disclosures, validated with reliability tests, and used as an auxiliary risk indicator in valuation and disclosure.

The RMB 203.7 billion estimate provides a benchmark for the firm-side cost of the 2021 regime. It is not a social cost, because it excludes consumer privacy benefits. It is also not the total value of data, because it covers only third-party data. It is, however, a useful order of magnitude for policy analysis: the adjustment cost is concentrated in a few industries and in a relatively small number of firms.

\subsection{Implications for Firms and Investors}

For firms, the results imply that third-party data dependence should be managed as a valuation risk. B2B data-service firms need to reassess the compliance cost of data procurement, client authorization, and cross-organizational sharing. Consumer firms with first-party data are less exposed, but they are not immune: first-party data did not fully buffer the shock in our sample. Building first-party data capability is valuable, but it is not a complete substitute for external data.

For investors, the TDDI provides an ex ante risk ranking. A text-based index can be computed before a regulatory shock, updated annually, and combined with industry and market-value data to estimate exposure. This is particularly relevant for data-intensive small and medium firms, whose data practices are less visible to analysts than those of large platforms.

\subsection{Scope of the Estimate}

The estimate applies to the 2021 privacy regime and to Chinese A-share digital firms. It captures the combined effect of ATT and PIPL, which overlap in time and operate through related margins. The TDDI is a text-based proxy for data dependence, and the sample covers an important segment of the digital economy. Within this scope, the estimated exposure is large and uneven across business models.
